%%%%%%%%%%%%%%%%%%%%%%%%%%%%%%%%%%%%%%%%%%%%%%%%%%%%%%%%%%%%%%%%%%%%%%
% Skupina20/main.tex -- Meja prepletenosti na procesorju IBM (stanja GHZ)
%
% Pravila predloge: brez \usepackage, slike v Skupina20/img/,
% vse oznake in ključi virov z "g20:", naši viri v references.bib v tej mapi
% (gradnja jih doda k references.bib predloge).
%
% Besedilo v oglatih oklepajih [ ... ] je še za dopolniti.
% Okvirje "Opomba za skupino" pred oddajo IZBRIŠEMO.
%%%%%%%%%%%%%%%%%%%%%%%%%%%%%%%%%%%%%%%%%%%%%%%%%%%%%%%%%%%%%%%%%%%%%%

\graphicspath{{Skupina20/img/}}

\title{Meja prepletenosti na procesorju IBM: koliko kubitov lahko zares prepletemo v stanje GHZ?}
\titlerunning{Meja prepletenosti na procesorju IBM}
\author{Dominik Krašovec in Erik Radovičevič}
\maketitle

\abstract{[Povzetek napišemo na koncu, ko imamo rezultate. Vsebovati mora: raziskovalno vprašanje (največje število kubitov $N_{\max}$, ki jih procesor IBM Heron prepleta v stanje GHZ z zvestobo $F>0{,}5$), metodo (populacije in paritetne oscilacije, linearna veriga proti drevesu, dinamično razklapljanje in ublažitev napak branja), glavni rezultat s številko, npr.\ \enquote{Linearna veriga doseže $N_{\max}=[\ ]$, drevo z ublažitvijo pa $N_{\max}=[\ ]$}, in en stavek o pomenu.]}


%======================================================================
\begin{svgraybox}
\textbf{Opomba za skupino (pred oddajo izbriši).} To poglavje je ogrodje, ki ga skupaj dopolnjujemo. Uvod in Metode so že
napisani tako, da iz njih razberete, \emph{kaj} bomo merili in \emph{zakaj}; Rezultati, Razprava in Zaključek so še prazni.
Pred začetkom preberite:
\begin{itemize}
\item razdelek~\ref{g20:sec:uvod} (kaj je stanje GHZ in zakaj je $F>0{,}5$ prava meja),
\item razdelek~\ref{g20:sec:merjenje} (kako $F$ izmerimo -- to je jedro naloge),
\item razdelek~\ref{g20:sec:cas} (procesorski čas je omejen na 600~s v 28 dneh, zato \textbf{nič ne poženemo na pravem procesorju brez dogovora}).
\end{itemize}
Koda je v \texttt{src/}, surovi rezultati z job ID v \texttt{rezultati/naloge.jsonl}. Roki: prijava teme 15.\,11.\,2026,
oddaja 22.\,11.\,2026, zagovor v tednu od 23.\,11.\,2026. Poročilo naj ima $\approx$\,10 strani.
\end{svgraybox}
%======================================================================


\section{Uvod}
\label{g20:sec:uvod}

Prepletenost je vir, ki kvantnim računalnikom omogoča prednost pred klasičnimi~\cite{nielsen2010}. Algoritmi, ki naj bi bili
uporabni, potrebujejo prepletenost čez veliko kubitov hkrati, današnji procesorji pa so šumni~\cite{preskill2018}: vsaka
dvokubitna vrata, vsak čas mirovanja in vsako branje prispevajo napako. Zato je naravno vprašanje, \emph{čez koliko kubitov
lahko danes dejansko ustvarimo prepletenost}. Na to vprašanje odgovarja priprava stanja Greenberger--Horne--Zeilinger
(GHZ)~\cite{g20:greenberger1990}
\begin{equation}
\ket{\mathrm{GHZ}_N} = \frac{1}{\sqrt{2}}\left(\ket{0}^{\otimes N} + \ket{1}^{\otimes N}\right),
\label{g20:eq:ghz}
\end{equation}
ki je \enquote{najbolj krhko} večdelno prepleteno stanje: ena sama napaka faze na kateremkoli kubitu uniči koherenco med
obema členoma. Stanje GHZ je zato strog in pogosto uporabljen preizkus kakovosti procesorja.

Mera, ki jo uporabimo, je zvestoba $F=\bra{\mathrm{GHZ}_N}\rho\ket{\mathrm{GHZ}_N}$ izmerjenega stanja $\rho$ glede na
idealno. Če velja $F>0{,}5$, je stanje dokazano \emph{pravo večdelno prepleteno} (angl.\ \emph{genuine multipartite
entanglement}), tj.\ ni ga mogoče zapisati kot mešanico stanj, v katerih bi bila katerakoli skupina kubitov ločena od
ostalih~\cite{g20:guhne2009}. Meja $0{,}5$ torej ni dogovor, ampak ima fizikalni pomen.

\textbf{Kaj so drugi že naredili.} Zvestobo stanj GHZ iz populacij in paritetnih oscilacij so prvič merili na ionih
(4~ioni~\cite{g20:sackett2000}, 14~ionov~\cite{g20:monz2011}). Na superprevodnih procesorjih IBM so poročali o 18-kubitnem
stanju z $F=0{,}517$~\cite{g20:wei2020} in o 27-kubitnem stanju z $F=0{,}546$ ob ublažitvi napak branja~\cite{g20:mooney2021}.
Leta 2025 je IBM z naprednimi tehnikami (prilagodljivo prevajanje, zaznavanje napak, začasno razračunavanje) pripravil
120-kubitno stanje GHZ z $F=0{,}56$~\cite{g20:ibm2025bigcats}. Pri predmetu NPMP v letih 2017/18--2025/26 nobena
seminarska naloga ni merila prepletenosti, ublažitve napak ali dinamičnega razklapljanja na realnem procesorju; skoraj vse so
implementacije znanih algoritmov.

\textbf{Kaj ostaja odprto.} Objavljeni rekordi uporabljajo tehnike in procesorski čas, ki jih študent z brezplačnim
paketom IBM Quantum Open Plan nima. Zanima nas, \emph{koliko daleč pridemo z vezji in ublažitvijo, ki so dostopni vsakomur}, in
kateri od teh vzvodov meji najbolj pomaga.

\begin{svgraybox}
\textbf{Raziskovalno vprašanje.} Koliko kubitov $N_{\max}$ lahko procesor IBM Heron (ibm\_fez, 156 kubitov) pripravi v stanje
GHZ z zvestobo $F>0{,}5$ in za koliko kubitov to mejo premakneta (a)~zasnova vezja (linearna veriga proti drevesu z majhno
globino) ter (b)~dinamično razklapljanje in ublažitev napak branja?

\textbf{Pričakovani rezultat.} Napoved klasičnega modela proračuna napak (razdelek~\ref{g20:sec:model}) iz kalibracije FakeFez:
največji $N$ z napovedanim $F>0{,}5$ je za linearno verigo $N_{\max}\approx 26$ brez ublažitve in $\approx 33$ z ublažitvijo
branja, za drevo, prilagojeno povezavam čipa, pa $\approx 33$ oziroma $\approx 70$. Surovo mejo določajo predvsem napake branja, mejo verige pa čakanje kubitov.
Dinamično razklapljanje bolj pomaga verigi, ker njeni kubiti dlje čakajo.
\end{svgraybox}

V razdelku~\ref{g20:sec:metode} opišemo vezja, način merjenja zvestobe in porabo procesorskega časa,
v razdelku~\ref{g20:sec:rezultati} rezultate na simulatorjih in procesorju, v razdelku~\ref{g20:sec:razprava} pa odgovorimo
na raziskovalno vprašanje in opišemo omejitve.


\section{Metode}
\label{g20:sec:metode}

\subsection{Orodja in nastavitve}
\label{g20:sec:orodja}

Vezja gradimo in izvajamo v okolju Qiskit~\cite{qiskit2024}, ker neposredno podpira procesorje IBM, njihove šumne modele in
dinamično razklapljanje. Različice in nastavitve so v tabeli~\ref{g20:tab:orodja}.

\begin{table}[htbp]
\caption{Orodja in nastavitve. [Dopolni ob glavnem zagonu: datum kalibracije ibm\_fez.]}
\label{g20:tab:orodja}
\centering
\begin{tabular}{ll}
\toprule
Parameter & Vrednost \\
\midrule
Qiskit & 2.6.1 \\
Qiskit Aer & 0.18.2 \\
qiskit-ibm-runtime & 0.51.0 \\
kvantni procesor & ibm\_fez (IBM Heron, 156 kubitov) \\
šumni simulator & Aer, metoda MPS, šumni model FakeFez (kalibracija 26.\,2.\,2025) \\
primitiv & SamplerV2 \\
prevajanje & \texttt{optimization\_level=1}, ročna postavitev kubitov, razporejanje ALAP \\
dinamično razklapljanje & X--X v času čakanja (\texttt{PadDynamicalDecoupling}) \\
ublažitev napak branja & lokalna, matrika $2\times2$ na kubit \\
faze za koherenco & $M=9$, $\phi_k=\pi k/9$ \\
število posnetkov (\emph{shots}) & 1000 na vezje, kalibracija $4\times1000$ na stanje \\
\bottomrule
\end{tabular}
\end{table}

Posle oddajamo s skripto \texttt{src/qpu.py}. Ukaz \texttt{nacrt} prevede vezja in brez oddaje izpiše njihove lastnosti,
oceno procesorskega časa po enačbi~\eqref{g20:eq:cas} in preostanek na instanci. Ukaz \texttt{poslji} zahteva zgornjo mejo
porabe (\texttt{-{}-max-sekund}), ki smo jo dogovorili pred zagonom. Vsakemu poslu nastavi največje trajanje
(\texttt{max\_execution\_time}) in ga pred oddajo v celoti rezervira v tej meji; rezervacijo sprosti šele, ko IBM posel obračuna.
Posel prekliče, če je ocena IBM večja od njegove meje, nov zagon pa zavrne, dokler kateri od prejšnjih poslov ni obračunan. Job ID, postavitev
kubitov in izide shrani v \texttt{rezultati/naloge.jsonl} in \texttt{rezultati/qpu/}.

\subsection{Vezja za pripravo stanja GHZ}
\label{g20:sec:vezja}

Stanje~\eqref{g20:eq:ghz} pripravimo z Hadamardovimi vrati na prvem kubitu in $N-1$ vrati CNOT, ki \enquote{širijo}
superpozicijo na ostale kubite. Število dvokubitnih vrat je pri obeh zasnovah enako ($N-1$), razlikuje pa se globina:

\begin{itemize}
\item \textbf{Linearna veriga} (slika~\ref{g20:fig:veriga}): CNOT-i si sledijo eden za drugim po poti kubitov na čipu.
      Globina raste kot $N$, zato prvi kubiti dolgo mirujejo in medtem dekoherirajo.
\item \textbf{Drevo} (slika~\ref{g20:fig:drevo}): vsak že prepleten kubit v naslednjem koraku prepleta svojega soseda, zato
      se število prepletenih kubitov v vsakem koraku lahko podvoji. Pri povezavah vsak-z-vsakim bi bila globina
      $\lceil\log_2 N\rceil$. Na procesorjih IBM so kubiti povezani v mrežo \emph{heavy-hex} (vsak kubit ima največ 3 sosede),
      zato drevo zgradimo z iskanjem v širino (BFS) iz korena z najmanjšo ceno (spodaj); globina je tedaj sorazmerna z
      razdaljo do najbolj oddaljenega kubita, približno $\sqrt{N}$.
\end{itemize}
Obe vezji gradimo neposredno na grafu povezav procesorja, tako da prevajalnik ne doda vrat SWAP. Kubite izberemo iz
kalibracije procesorja (napake branja $e_q$ in napake dvokubitnih vrat $e_{qr}$):
\begin{itemize}
\item \textbf{veriga} je najcenejša od enostavnih poti z $N$ kubiti, ki jih požrešno iskanje v globino z vračanjem najde iz vsakega
      začetnega kubita; cena poti je
      $\sum_q e_q + \sum_{(q,r)} e_{qr}$;
\item \textbf{drevo} zgradimo z BFS iz vsakega možnega korena in obdržimo tisto z najmanjšo ceno
      $\sum_q e_q + \sum_{(q,r)} e_{qr} + 0{,}0015\,N h$, kjer je $h$ globina drevesa. Zadnji člen grobo oceni izgubo koherence
      med čakanjem: vsak nivo drevesa pomeni $\approx 0{,}15~\mu$s čakanja za $\approx N$ kubitov, kar pri $T_2\approx100~\mu$s
      da $\approx 0{,}0015$ na kubit in nivo. Prva različica je za koren vzela osrednji kubit in je pri $N=50$ izbrala
      kubite s slabšim branjem od verige (povprečna napaka branja 0,0132 proti 0,0100).
\end{itemize}
Za vsako vezje poročamo globino in število dvokubitnih vrat pred prevajanjem in po njem ter trajanje prevedenega vezja.

\begin{figure}[htbp]
\centering
\begin{subfigure}[b]{0.48\textwidth}
\centering
\begin{quantikz}[row sep=0.25cm, column sep=0.25cm]
\lstick{$q_0$} & \gate{H} & \ctrl{1} & \qw      & \qw      & \qw \\
\lstick{$q_1$} & \qw      & \targ{}  & \ctrl{1} & \qw      & \qw \\
\lstick{$q_2$} & \qw      & \qw      & \targ{}  & \ctrl{1} & \qw \\
\lstick{$q_3$} & \qw      & \qw      & \qw      & \targ{}  & \qw
\end{quantikz}
\caption{Linearna veriga, globina $N$.}
\label{g20:fig:veriga}
\end{subfigure}
\hfill
\begin{subfigure}[b]{0.48\textwidth}
\centering
\begin{quantikz}[row sep=0.25cm, column sep=0.25cm]
\lstick{$q_2$} & \qw      & \qw      & \targ{}   & \qw \\
\lstick{$q_0$} & \gate{H} & \ctrl{1} & \ctrl{-1} & \qw \\
\lstick{$q_1$} & \qw      & \targ{}  & \ctrl{1}  & \qw \\
\lstick{$q_3$} & \qw      & \qw      & \targ{}   & \qw
\end{quantikz}
\caption{Drevo, globina $\approx\log_2 N$ pri povezavah vsak-z-vsakim.}
\label{g20:fig:drevo}
\end{subfigure}
\caption{Priprava 4-kubitnega stanja GHZ z dvema zasnovama vezja. V drevesu se zadnja CNOT-a izvedeta hkrati, zato je
vezje za en korak plitvejše; pri večjih $N$ razlika hitro raste. [Kasneje dodamo izris dejanske postavitve kubitov na ibm\_fez.]}
\label{g20:fig:vezja}
\end{figure}

\subsection{Merjenje zvestobe}
\label{g20:sec:merjenje}

Zvestobe ne moremo izmeriti z eno samo meritvijo. Za stanje GHZ jo razstavimo na dva dela, ki ju izmerimo ločeno~\cite{g20:sackett2000, g20:mooney2021}:
\begin{equation}
F = \frac{P + C}{2}.
\label{g20:eq:zvestoba}
\end{equation}

\textbf{Populacija $P$.} Vse kubite izmerimo v računski bazi; $P=p(0\dots0)+p(1\dots1)$ je delež posnetkov, ki dajo enega
od obeh pravilnih izidov. Idealno je $P=1$. Visok $P$ še ne pomeni prepletenosti -- enak $P$ bi dala tudi klasična mešanica
\enquote{vse 0 ali vse 1}.

\textbf{Koherenca $C$ (paritetne oscilacije).} Pred meritvijo na vsak kubit apliciramo rotacijo, ki ga meri v bazi
$\cos\phi\,X+\sin\phi\,Y$, in izračunamo pariteto $\Pi=\expval{Z^{\otimes N}}$ (delež sodih minus delež lihih izidov). Za
stanje GHZ velja
\begin{equation}
\expval{\Pi(\phi)} = C\cos(N\phi + \phi_0),
\label{g20:eq:pariteta}
\end{equation}
kjer je amplituda $C$ ravno velikost koherence med $\ket{0\dots0}$ in $\ket{1\dots1}$; idealno $C=1$, pri klasični mešanici
$C=0$. Zato klasična mešanica da $F\le 0{,}5$ in meje ne more preseči.

Pri šumnem stanju enačba~\eqref{g20:eq:pariteta} ne velja natančno. Pariteta je tedaj vsota členov s frekvencami $N-2j$,
$j=0,\dots,N$: člen z $j>0$ izvira iz koherence med komplementarnima nizoma bitov, ki se od $0\dots0$ in $1\dots1$ razlikujeta v
$j$ bitih, in ga povzroči šum. Signal ločimo z diskretno Fourierovo transformacijo pri $M$ fazah na polovici kroga,
\begin{equation}
c_N = \frac{1}{M}\sum_{k=0}^{M-1} \expval{\Pi(\phi_k)}\, e^{-iN\phi_k}, \qquad \phi_k=\frac{\pi k}{M}, \qquad C = 2\abs{c_N}.
\label{g20:eq:dft}
\end{equation}
Vsota $\sum_k e^{-i2j\phi_k}$ je enaka nič za vsak $j$, ki ni večkratnik $M$, zato ostanejo le členi z $j=M, 2M,\dots$, ki
zahtevajo $M$ hkratnih preklopov in so zanemarljivi. Člen z $j=N$ je konjugirani signal (frekvenca $-N$), zato $N$ ne sme
biti večkratnik $M$. Uporabimo $M=9$, kar da 10 vezij na točko (1 populacijsko in 9 paritetnih) namesto $2N+3$ pri polnem
naboru faz $\phi_j=\pi j/(N+1)$, $j=0,\dots,2N+1$~\cite{g20:wei2020}. V šumni simulaciji FakeFez se $C$ iz 8 faz od polnega
nabora razlikuje za 0,009 pri $N=10$ in 0,004 pri $N=13$, kar je znotraj negotovosti. [Izid shrani v \texttt{rezultati/sim/}.]

\textbf{Negotovost in kriterij.} Negotovost $F$ ocenimo z zankanjem (angl.\ \emph{bootstrap}, 200 ponovitev): števila izidov
vsakega vezja, tudi kalibracijskih, na novo vzorčimo iz multinomske porazdelitve in $F$ izračunamo znova, zato je negotovost
popravka branja del $\sigma_F$. Za vsako konfiguracijo je $N_{\max}$ največji $N$, pri katerem je $F-2\sigma_F>0{,}5$. Ker
merimo le nekaj vrednosti $N$, $N_{\max}$ podamo kot interval med največjim $N$, ki kriterij izpolni, in najmanjšim, ki ga ne.

\subsection{Ublažitev napak}
\label{g20:sec:ublazitev}

\textbf{Dinamično razklapljanje (DD).} V času, ko kubit čaka, nanj apliciramo zaporedje impulzov, katerih skupni učinek je
identiteta, vendar sproti izničijo počasno fazno napako zaradi okolice~\cite{g20:viola1999, g20:ezzell2023}. Vezje ob
prevajanju razporedimo ALAP (\enquote{čim pozneje}), tako da čakanja postanejo eksplicitni ukazi \texttt{delay}, in v vsako
čakanje vstavimo par impulzov X (\texttt{PadDynamicalDecoupling} v Qiskitu; X--X je tudi privzeto zaporedje IBM). Vezji brez
DD in z DD imata enaka čakanja in sta v istem poslu, zato ju ne loči sprememba kalibracije med posli. DD naj bi posebej
pomagalo linearni verigi, kjer kubiti dolgo mirujejo. Šumni simulator pozna le markovsko relaksacijo $T_1$, $T_2$, proti kateri
DD ne pomaga, zato učinek DD izmerimo samo na procesorju.

\textbf{Ublažitev napak branja.} Kalibracijski vezji pripravita vse kubite v $\ket{0}$ oziroma $\ket{1}$ in ju takoj izmerita.
Iz njiju za vsak kubit $i$ dobimo $e_{0,i}=p(1\,|\,0)$ in $e_{1,i}=p(0\,|\,1)$ ter matriko
\begin{equation}
A_i=\begin{pmatrix}1-e_{0,i} & e_{1,i}\\ e_{0,i} & 1-e_{1,i}\end{pmatrix},
\label{g20:eq:branje}
\end{equation}
ki preslika pravo stanje kubita v prebrano. Napake branja različnih kubitov obravnavamo kot neodvisne (lokalna ublažitev,
kot v~\cite{g20:mooney2021}). Popravljena vrednost $p(0\dots0)$, $p(1\dots1)$ ali paritete je tedaj utežena vsota izidov, kjer
je utež izida produkt elementov $A_i^{-1}$ po kubitih, zato matrike velikosti $2^N\times2^N$ ne potrebujemo. To je naknadna
obdelava; procesorski čas porabita le kalibracijski vezji. Vsak posel ju vsebuje enkrat za vse kubite, ki jih uporablja, vsako
ponovljeno $4\times$ po 1000 posnetkov, ker morajo imeti vsa vezja v poslu enako število posnetkov.

\subsection{Šumni simulator in klasični model}
\label{g20:sec:model}

\textbf{Šumni simulator.} Simuliramo z Aer po metodi matrično-produktnih stanj (MPS); stanje GHZ ima vezno dimenzijo 2, zato je
simulacija tudi pri 50 kubitih hitra. Šumni model vzamemo iz FakeFez (posnetek kalibracije ibm\_fez s 26.\,2.\,2025). Ker so
vezja razporejena ALAP, model na ukazih \texttt{delay} uporabi relaksacijo $T_1$, $T_2$ (Bellov par po 20~$\mu$s čakanja: pariteta
$XX$ pade z 0,98 na 0,66 [preverbo shrani v \texttt{rezultati/sim/}]). Od modela odstopamo na dveh mestih:
\begin{itemize}
\item napako dvokubitnih vrat nadomestimo z njenim Paulijevim približkom (verjetnosti $p_P=\sum_k\abs{\tr(P^\dagger K_k)}^2/16$
      za Krausove operatorje $K_k$), ker simulacija MPS pri dvokubitnih kanalih Krausa obtiči
      (razdelek~\ref{g20:sec:neuspehi}); izgubi se le neunitalni del napake med samimi vrati;
\item napake branja dodamo po simulaciji kot klasične preklope bitov z nesimetričnima verjetnostma $p(1\,|\,0)$ in
      $p(0\,|\,1)$ iz kalibracije, ker Aerov model hrani le njuno povprečje.
\end{itemize}

\textbf{Klasični model proračuna napak.} Za vsako prevedeno in razporejeno vezje izračunamo verjetnost, da nobena vrata ne
naredijo napake, $G=\prod_g(1-e_g)$, ter razpad in izgubo koherence med čakanji,
\begin{equation}
A_1=\prod_i e^{-\tau_i/T_{1,i}}, \qquad D=\prod_i e^{-\tau_i/T_{2,i}},
\label{g20:eq:model}
\end{equation}
kjer je $\tau_i$ skupni čas čakanja kubita $i$ po njegovih prvih vratih , pri čemer $T_2$ omejimo na $2T_1$. Z idealno ublažitvijo branja je $P=G(1+A_1)/2$ in
$C=GD$, brez ublažitve pa $P=G\,[\prod_i(1-e_{0,i})+A_1\prod_i(1-e_{1,i})]/2$ in $C=GD\prod_i(1-e_{0,i}-e_{1,i})$; v obeh
primerih $F=(P+C)/2$. Model je grob (vsaka napaka vrat uniči stanje, čakanja so neodvisna), a ne potrebuje simulacije, zato
napove $F$ za vse $N$. Je naša klasična primerjava in osnova za izbiro $N$ v poskusih na procesorju.

\subsection{Zasnova poskusov}
\label{g20:sec:poskusi}

\begin{enumerate}
\item \textbf{Idealni simulator} (Aer, MPS): preverimo, da vezja in analiza dajo $F=1$ za obe zasnovi (samodejni testi v
      \texttt{tests/}).
\item \textbf{Ponovitev prejšnjega dela}: naš 5-kubitni zagon na ibm\_fez z 2.\,11.\,2026 ponovimo na šumnem simulatorju
      FakeFez z enakim vezjem in številom posnetkov ter primerjamo $P$.
\item \textbf{Šumni simulator} FakeFez pri majhnih $N$ in klasični model za vse $N$: preverba modela in napoved $N_{\max}$.
\item \textbf{Pilotni zagon} na ibm\_fez: 1~posel, 88~vezij; veriga $N=10, 22$ in drevo $N=20, 40$, vsaka točka brez in z
      DD. Preverimo model porabe časa; če je zaračunani čas na posnetek več kot 1,3-krat večji od modela, ustavimo.
\item \textbf{Glavni zagon}: 1~posel, 148~vezij; veriga $N=16, 28, 34$ in drevo $N=30, 50, 64, 76$, vsaka točka brez in z DD.
      Seznama $N$ po pilotu premaknemo, če ibm\_fez očitno odstopa od modela.
\item Po potrebi \textbf{zgostitev}: 1~posel z do dvema dodatnima točkama okoli meje $N_{\max}$.
\end{enumerate}
Vsaka točka (zasnova, $N$, DD) ima 10 vezij, posel pa še 8 kalibracijskih ($2\times4$). Ker strošek točke ni odvisen od $N$,
lahko verigo merimo okoli njene meje, drevo pa dlje.

\subsection{Poraba procesorskega časa}
\label{g20:sec:cas}

Paket Open Plan omogoča 600~s procesorskega časa v drsečem 28-dnevnem oknu. IBM porabo ocenjuje kot~\cite{g20:ibmjobtime}
\begin{equation}
t \approx 2\,\mathrm{s}\cdot n_\text{poslov} + 0{,}35\,\mathrm{ms}\cdot n_\text{vezij}\cdot n_\text{posnetkov},
\label{g20:eq:cas}
\end{equation}
kar se ujema s prvim zagonom (1 vezje $\times$ 1000 posnetkov: ocena 2,27~s, zaračunano 2~s). Ker je fiksni strošek posla
velik, vezja združujemo v malo poslov. Načrtovana poraba je v tabeli~\ref{g20:tab:cas}.

\begin{table}[htbp]
\caption{Načrtovana poraba procesorskega časa po enačbi~\eqref{g20:eq:cas}. [Dopolni z dejansko porabo, \texttt{job.usage()}.]}
\label{g20:tab:cas}
\centering
\begin{tabular}{lrrrr}
\toprule
Del poskusa & vezja & posli & ocena [s] & dejansko [s] \\
\midrule
prvi preizkus (5 kubitov) & 1 & 1 & 2,3 & 2 \\
pilotni zagon & 88 & 1 & 32,8 & [ ] \\
glavni zagon & 148 & 1 & 53,8 & [ ] \\
zgostitev (po potrebi) & $\le 48$ & 1 & $\le 18{,}8$ & [ ] \\
\midrule
skupaj & $\le 285$ & $\le 4$ & $\le 107{,}7$ & [ ] \\
\bottomrule
\end{tabular}
\end{table}

\subsection{Metrike}
\label{g20:sec:metrike}

\begin{itemize}
\item zvestoba $F$, populacija $P$ in koherenca $C$ v odvisnosti od $N$ (enačbi~\eqref{g20:eq:zvestoba} in~\eqref{g20:eq:pariteta}),
\item $F$ surovo in z ublažitvijo branja, $N_{\max}$ za vsako konfiguracijo in prirastek $\Delta N_{\max}$ zaradi zasnove
      vezja, DD oziroma ublažitve branja,
\item razlika $F$ med procesorjem, šumnim simulatorjem in klasičnim modelom,
\item globina vezja in število dvokubitnih vrat pred prevajanjem in po njem, trajanje prevedenega vezja,
\item porabljen procesorski čas na konfiguracijo.
\end{itemize}


\section{Rezultati}
\label{g20:sec:rezultati}

\textbf{Prvi preizkus na procesorju.} 2.\,11.\,2026 smo na ibm\_fez pognali 5-kubitno linearno verigo
(1000 posnetkov, job ID \texttt{dbk48fl6p51s73clupeg}). Po prevajanju je imelo vezje globino 20 in 4 vrata CZ. Porazdelitev
izidov je na sliki~\ref{g20:fig:ghz5}: $P = 0{,}460 + 0{,}395 = 0{,}855$. Asimetrija med $\ket{00000}$ in $\ket{11111}$ je
pričakovana, ker se kubit v stanju $\ket{1}$ med čakanjem na branje lahko sprosti v $\ket{0}$. Koherence $C$ v tem preizkusu
še nismo merili, zato zvestobe še ne poznamo.

\begin{figure}[htbp]
\centering
\includegraphics[width=0.85\textwidth]{ghz5_ibm}
\caption{Porazdelitev izidov 5-kubitnega stanja GHZ na procesorju ibm\_fez (1000 posnetkov). Izidi z manj kot 5 zadetki so
združeni v stolpec \enquote{ostali}. Sliko izriše \texttt{src/izris\_ghz5.py}.}
\label{g20:fig:ghz5}
\end{figure}

\textbf{Ponovitev na šumnem simulatorju.} Fizični kubiti zagona z 2.\,11.\,2026 niso bili zabeleženi, zato smo enako vezje
(1000 posnetkov) simulirali na najboljši verigi FakeFez (kubiti 141, 142, 143, 136, 123) in na 10 naključnih verigah
(\texttt{rezultati/sim/ponovitev\_B0.json}). Prevedeno vezje ima tudi tu globino 20 in 4 vrata CZ. Najboljša veriga da
$P=0{,}948\pm0{,}007$, naključne pa $P=0{,}842$--$0{,}923$; ibm\_fez s $P=0{,}855$ leži na spodnjem robu tega razpona.
[Razlaga: postavitev kubitov na ibm\_fez ni znana, kalibracija FakeFez je s 24.\,2.\,2025.]

\textbf{Klasični model proti šumni simulaciji.} Tabela~\ref{g20:tab:model} primerja napoved modela
(razdelek~\ref{g20:sec:model}) s šumno simulacijo verige na FakeFez. Razlika je največ 0,016; model je v petih od šestih
primerov nekoliko višji od simulacije. [Razlaga odstopanja: \ ] [Točki $N=10$ in 13 shrani v \texttt{rezultati/sim/}.]

\begin{table}[htbp]
\caption{Zvestoba verige na FakeFez: klasični model proti šumni simulaciji (MPS).}
\label{g20:tab:model}
\centering
\begin{tabular}{rrrrr}
\toprule
& \multicolumn{2}{c}{surovo} & \multicolumn{2}{c}{z ublažitvijo branja} \\
$N$ & model & simulacija & model & simulacija \\
\midrule
5 & 0,944 & 0,932 & 0,982 & 0,972 \\
10 & 0,854 & 0,850 & 0,942 & 0,945 \\
13 & 0,777 & 0,764 & 0,887 & 0,871 \\
\bottomrule
\end{tabular}
\end{table}

[Tukaj pridejo: (1) tabela lastnosti vezij za vse $N$ in obe zasnovi; (2) graf $F(N)$ za idealni simulator, FakeFez,
klasični model in ibm\_fez, vse 4 konfiguracije, z vodoravno črto pri $F=0{,}5$; (3) tabela $N_{\max}$ po konfiguracijah;
(4) primerjava z objavami~\cite{g20:wei2020, g20:mooney2021, g20:ibm2025bigcats}.]


\section{Razprava}
\label{g20:sec:razprava}

[Odgovor na raziskovalno vprašanje: ali se je hipoteza potrdila? Kateri vzvod (zasnova vezja, DD, ublažitev branja) je meji
najbolj pomagal in zakaj? Kako dobro je šumni model FakeFez napovedal rezultat procesorja?]

\textbf{Omejitve.} [Npr.: kalibracija procesorja se med dnevi spreminja; omejen procesorski čas dovoljuje le eno
ponovitev; ublažitev branja popravi branje, ne pa napak med pripravo stanja.]

\subsection{Kaj ni delovalo}
\label{g20:sec:neuspehi}

[Opomba za skupino: zahtevan je vsaj en neuspešen poskus; izberemo enega ali dva od spodnjih in dodamo morebitne s procesorja.]

\textbf{Simulacija MPS je obtičala.} Simulacija 27-kubitnega stanja GHZ s šumnim modelom FakeMontrealV2 se ni nikoli končala
in ni javila napake. Prožilo jo je termično sproščanje v dvokubitnem kanalu Krausa pri vratih CX, pojavljala pa se je šele pri
več posnetkih (pri 200 posnetkih je delovala, pri 4000 ne). Rešitev sta Paulijev približek napake dvokubitnih vrat in
dodajanje napak branja po simulaciji (razdelek~\ref{g20:sec:model}). Na FakeFez je približek 5-kubitni $F$ spremenil le z
0,978 na 0,974 [izid ni shranjen: ponovi ali številki izbriši].

\textbf{Simetrične napake branja.} Aerov šumni model napako branja kubita hrani kot povprečje $p(1\,|\,0)$ in $p(0\,|\,1)$,
zato simulacija ni kazala asimetrije med $\ket{0\dots0}$ in $\ket{1\dots1}$, ki smo jo videli na ibm\_fez (460 proti 395
zadetkov). Zdaj obe verjetnosti beremo iz kalibracije procesorja vsako posebej.

\textbf{Ponovitev objav na simulatorjih starih procesorjev.} Načrtovali smo tudi ponovitev rezultatov~\cite{g20:wei2020, g20:mooney2021} na šumnih modelih takratnih procesorjev (FakeMontrealV2 in 20-kubitni modeli). Dolge simulacije na vseh jedrih so prenosnik večkrat pregrele do izklopa. Ponovitvi smo opustili, ker za raziskovalno vprašanje nista potrebni: z objavami neposredno primerjamo naše točke na ibm\_fez pri $N=20$--$30$.

\textbf{Poraba 0~s.} Naša prva skripta za preverjanje dostopa je po prvem zagonu izpisala porabo 0~s, ker je porabo prebrala,
preden jo je IBM obračunal; dejanska poraba je bila 2~s. Enaka past bi pokvarila tudi proračun: če bi porabo šteli sproti,
bi lahko oddali več, kot smo dogovorili. Zato \texttt{src/qpu.py} vsakemu poslu vnaprej rezervira njegovo največje trajanje in
rezervacijo sprosti šele po obračunu (razdelek~\ref{g20:sec:orodja}).


\section{Zaključek}
\label{g20:sec:zakljucek}

[Glavni rezultat v nekaj stavkih in predlog nadaljnjega dela, npr.\ dinamična vezja z meritvami sredi vezja, ki omogočajo
pripravo GHZ v konstantni globini.]


\section*{Doprinosi avtorjev}

[Za vsakega člana (Dominik, Erik) dva do tri stavke. Dogovor: oba delava na vseh delih (vezja in izbira kubitov, simulacije
in zagoni na IBM, analiza, literatura, uvod, razprava, prosojnice, README); vsak del napiše ali izvede eden, drugi ga pregleda
in neodvisno ponovi. Tu iskreno napiši, kdo je kaj dejansko naredil.]

\section*{Izjava o uporabi generativne umetne inteligence}

Pri razvoju smo si pomagali z orodjem generativne umetne inteligence: pri iskanju informacij, razlagi konceptov, pisanju kode,
pripravi in analizi poskusov ter pisanju in urejanju besedila poročila.
